\paragraph{扫描剖面的标量证书：混叠阈值、稳定反演与范数夹逼}
在有限缺陷口径下，共动扫描剖面不仅提供可用于 Prony 反演的指数谱指纹，还同时诱导若干仅依赖标量统计量的可审计证书。
以下给出在制度语义中最直接的混叠阈值、稳定性与范数下界。

\begin{proposition}[地址离散化的混叠阈值：格点采样仅确定 \texorpdfstring{$\gamma$}{gamma} 的同余类]\label{prop:xi-scan-aliasing-threshold}
\leanverified{paper\_xi\_scan\_aliasing\_threshold}
固定采样步长 \(\Delta>0\)，并仅观测谱指纹离散样本 \(\{u_n\}_{n\ge 0}\)（\(u_n=\widehat H(n\Delta)\)）。
则对任意整数 \(k\in\mathbb{Z}\) 与任意 \(j\)，把高度参数平移
\(
\gamma_j\mapsto \gamma_j+\frac{2\pi k}{\Delta}
\)
不会改变任何样本 \(u_n\)。
因此离散采样只能确定高度参数的同余类 \(\gamma_j\ \mathrm{mod}\ (2\pi/\Delta)\)。
\par
若先验高度窗为 \(|\gamma_j|\le T\)，则实现唯一 lifting（无混叠）的必要条件为
\[
\boxed{\;\Delta<\frac{\pi}{T}.\;}
\]
\end{proposition}

\begin{proof}
由 \eqref{eq:xi-comoving-scan-fourier}，对 \(n\ge 0\) 有
\[
u_n
:=4\pi\sum_{j=1}^{\kappa}m_j\,\frac{\delta_j}{1+\delta_j}\,e^{-(1+\delta_j)n\Delta}\,e^{-\mathrm{i}\gamma_j n\Delta}.
\]
若 \(\gamma_j\) 替换为 \(\gamma_j+\frac{2\pi k}{\Delta}\)，则
\(
e^{-\mathrm{i}(\gamma_j+2\pi k/\Delta)n\Delta}
:=e^{-\mathrm{i}\gamma_j n\Delta}e^{-2\pi\mathrm{i}kn}
:=e^{-\mathrm{i}\gamma_j n\Delta}
\)，
故所有 \(u_n\) 不变，第一部分成立。
\par
若 \(\Delta\ge \pi/T\)，则 \(2\pi/\Delta\le 2T\)，取 \(\gamma=-T\) 则 \(\gamma':=\gamma+2\pi/\Delta\in[-T,T]\)，
但二者在离散样本上不可区分，从而 lifting 非唯一；故无混叠必要求 \(\Delta<\pi/T\)。证毕。
\end{proof}

\begin{theorem}[有限采样的最小完备性：\texorpdfstring{$2\kappa$}{2kappa} 阶谱指纹的必要且一般充分性]\label{thm:xi-scan-minimal-complete-2kappa}
在定理 \ref{thm:xi-comoving-scan-hankel-rank-defect} 的设定下，设谱点 \(z_j\) 两两不同。
则由连续样本
\[
\boxed{\;
(u_0,u_1,\dots,u_{2\kappa-1})\;}
\]
可唯一恢复 \(\{(z_j,A_j)\}_{j=1}^{\kappa}\)（忽略排列），从而结合定理 \ref{thm:xi-comoving-scan-exponential-spectral-separation} 的闭式反演唯一确定
\(\{(\gamma_j,\delta_j,m_j)\}_{j=1}^{\kappa}\)。
\par
反之，若仅给定长度为 \(N+1\) 的前缀样本 \((u_0,\dots,u_{N})\) 且 \(N\le 2\kappa-2\)，
则存在两组不同的有限缺陷参数（同一缺陷度 \(\kappa\)）产生相同的该长度样本序列；
因此在一般位置假设下，\(2\kappa\) 为最小完备采样长度。
\end{theorem}

\begin{proof}
由 \(\mathrm{rank}(\mathsf H^{(\kappa-1)})=\kappa\)（定理 \ref{thm:xi-comoving-scan-hankel-rank-defect}），Hankel 线性系统的 annihilating 多项式（Prony 多项式）在一般位置下由阶 \(\kappa\) 的核关系唯一确定，从而 \(\{z_j\}\)（作为该多项式的根）可由 \(u_0,\dots,u_{2\kappa-1}\) 唯一恢复；随后 \(A_j\) 由 Vandermonde 系统唯一求出。
因此 \(\{(z_j,A_j)\}\) 唯一，进而由 \(|z_j|\) 与 \(\arg(z_j)\) 分离出 \((\delta_j,\gamma_j)\)，再由 \eqref{eq:xi-comoving-scan-fourier} 的振幅关系恢复 \(m_j\)。
\par
最小性可由参数维度与观测维度的欠定性给出：映射
\(
\{(z_j,A_j)\}_{j=1}^{\kappa}\mapsto (u_0,\dots,u_N)
\)
在 \(N+1<2\kappa\) 时一般欠定（未知量维度严格大于观测维度），故其纤维一般具有正维数，从而存在不同参数产生相同前缀样本。
这表明 \(2\kappa\) 阶样本在一般位置下不可再压缩。证毕。
\leanverified{paper\_xi\_scan\_minimal\_complete\_2kappa}
\end{proof}

\begin{theorem}[最小分离 \texorpdfstring{$\Rightarrow$}{=>} 反演稳定性：层析反演在几何间距约束下局部双 Lipschitz]\label{thm:xi-scan-min-separation-local-bilipschitz}
\leanverified{paper\_xi\_scan\_min\_separation\_local\_bilipschitz}
在定理 \ref{thm:xi-comoving-scan-exponential-spectral-separation} 的无混叠口径下，
若存在分离常数 \(s>0\) 使得 \(|\gamma_j-\gamma_k|\ge s\)（\(j\ne k\)），并且 \(\delta_j\in[\delta_{\min},\delta_{\max}]\subset(0,\infty)\)，
则参数到有限谱指纹的映射
\[
\{(\gamma_j,\delta_j,m_j)\}_{j=1}^{\kappa}\ \longmapsto\ \{u_n\}_{n=0}^{2\kappa-1}
\]
在该受限参数流形上为局部双 Lipschitz：存在常数 \(C_1,C_2>0\)（仅依赖 \(\kappa,s,\delta_{\min},\delta_{\max}\) 与制度固定的 \(\Delta\)）使得对任意两组满足同分离约束的参数 \(\Theta,\widetilde\Theta\) 有
\[
C_1\,d(\Theta,\widetilde\Theta)
\ \le\
\max_{0\le n\le 2\kappa-1}|u_n(\Theta)-u_n(\widetilde\Theta)|
\ \le\
C_2\,d(\Theta,\widetilde\Theta),
\]
其中 \(d\) 可取为对称化的最优配对距离（忽略点的排列）：
\[
d(\Theta,\widetilde\Theta)
:=
\min_{\sigma\in S_\kappa}\max_{1\le j\le \kappa}\bigl(|\gamma_j-\widetilde\gamma_{\sigma(j)}|+|\delta_j-\widetilde\delta_{\sigma(j)}|\bigr).
\]
（在本文制度语义中 \(m_j\) 为重数数据并可视作离散已知；若允许权重连续扰动，则上述距离需额外补入权重项。）
\end{theorem}

\begin{proof}
由 \eqref{eq:xi-comoving-scan-fourier}，离散样本满足指数和表示 \(u_n=\sum_{j=1}^{\kappa}A_j z_j^{\,n}\)，其中
\(
z_j=e^{-(1+\delta_j)\Delta}e^{-\mathrm{i}\gamma_j\Delta}
\)
且
\(
A_j=4\pi m_j\frac{\delta_j}{1+\delta_j}>0
\)。
在分离与有界偏移约束下，节点集 \(\{z_j\}\) 落在紧集内且两两分离，从而相应的 Vandermonde/Prony 线性系统在局部可逆，并且其条件数在该紧集上有一致界。
因此由解析依赖与逆函数定理可得局部双 Lipschitz；常数依赖性由上述一致条件数界给出。证毕。
\end{proof}

\begin{corollary}[碰撞不适定性：缺陷合并导致反演条件数的必然爆炸]\label{cor:xi-scan-collision-illposedness}
在定理 \ref{thm:xi-scan-minimal-complete-2kappa} 的一般位置设定下，记最小谱间距
\[
\boxed{\;
\mathrm{sep}:=\min_{j\neq k}|z_j-z_k|\;}
\]
并以 \(L\) 表示任一基于有限样本 \((u_n)_{n=0}^{2\kappa-1}\) 的局部反演（包括 Prony/Hankel 核空间法及其等价变体）在该点处的 Lipschitz 常数（误差放大因子）。
则存在常数 \(C_\kappa>0\)（仅依赖 \(\kappa\) 与制度固定规范）使得
\[
\boxed{\;
L\ \ge\ C_\kappa\,\mathrm{sep}^{-(\kappa-1)}.\;}
\]
特别地，当两缺陷满足 \(\delta_j=\delta_k\) 且 \(|\gamma_j-\gamma_k|\to 0\) 时，有 \(\mathrm{sep}\asymp \Delta|\gamma_j-\gamma_k|\)，从而任何稳定反演的误差放大因子至少按 \(|\gamma_j-\gamma_k|^{-(\kappa-1)}\) 发散。
\leanverified{paper\_xi\_scan\_collision\_illposedness}
\end{corollary}

\begin{proof}
在一般位置下，\(\mathsf H^{(\kappa-1)}\) 满秩且具有因子分解
\(
\mathsf H^{(\kappa-1)}=V\,\mathrm{diag}(A_1,\dots,A_\kappa)\,V^{\mathsf T}
\)，其中 \(V\) 为 \(\kappa\times\kappa\) 的 Vandermonde 矩阵（见定理 \ref{thm:xi-comoving-scan-hankel-rank-defect} 的证明）。
因此
\[
\det\mathsf H^{(\kappa-1)}=(\det V)^2\prod_{j=1}^{\kappa}A_j,
\qquad
|\det V|=\prod_{1\le j<k\le \kappa}|z_k-z_j|
\ \le\ \mathrm{sep}^{\kappa(\kappa-1)/2}.
\]
任一等价反演都需在有限步内解出相应的 Hankel/Prony 线性系统，其稳定性下界由 \(\|(\mathsf H^{(\kappa-1)})^{-1}\|\) 控制；
而 \(\|(\mathsf H^{(\kappa-1)})^{-1}\|\) 至少与 \(\sigma_{\min}(\mathsf H^{(\kappa-1)})^{-1}\) 同阶，并可由行列式与谱范数之间的基本不等式下界估计，从而当 \(\mathrm{sep}\downarrow 0\) 时误差放大因子必以至少 \(\mathrm{sep}^{-(\kappa-1)}\) 的幂次发散（该幂次来自 \((\det V)^2\) 的平方 Vandermonde 因子）。
最后一条由 \(z_j=e^{-(1+\delta)\Delta}e^{-\mathrm{i}\gamma_j\Delta}\) 的一阶展开给出：当 \(\delta\) 固定且 \(|\gamma_j-\gamma_k|\) 很小时，\(|e^{-\mathrm{i}\gamma_j\Delta}-e^{-\mathrm{i}\gamma_k\Delta}|\asymp \Delta|\gamma_j-\gamma_k|\)，故 \(\mathrm{sep}\asymp \Delta|\gamma_j-\gamma_k|\)。证毕。
\end{proof}

\begin{proposition}[Cauchy--Poisson 剖面的 \texorpdfstring{$L^1$--$L^2$--$L^\infty$}{L1-L2-Linfty} 统一下界：缺陷必产生可见峰值]\label{prop:xi-scan-profile-lp-peak-witness}
令 \(H\) 如 \eqref{eq:xi-comoving-scan-profile-H} 所定义（故 \(H\ge 0\)）。
则严格恒等式成立
\[
\boxed{\;
\|H\|_{L^1(\RR)}
:=
4\pi\sum_{j=1}^{\kappa}m_j\,\frac{\delta_j}{1+\delta_j}\;}
\]
并且有平方积分下界
\[
\boxed{\;
\|H\|_{L^2(\RR)}^2
\ \ge\
\sum_{j=1}^{\kappa}8\pi\,m_j^{2}\,\frac{\delta_j^{2}}{(1+\delta_j)^{3}}\;}
\]
从而存在 \(x^\ast\in\RR\) 使
\[
\boxed{\;
\|H\|_{L^\infty(\RR)}
\ \ge\
\frac{\|H\|_{L^2(\RR)}^2}{\|H\|_{L^1(\RR)}}
\ \ge\
2\cdot
\frac{\sum_{j=1}^{\kappa} m_j^2\,\dfrac{\delta_j^2}{(1+\delta_j)^3}}
{\sum_{j=1}^{\kappa} m_j\,\dfrac{\delta_j}{1+\delta_j}}\;}
\]
即只要点状缺陷非空，扫描剖面必出现严格可检测的峰值下界。
\leanverified{paper\_xi\_scan\_profile\_lp\_peak\_witness}
\end{proposition}

\begin{proof}
对每一峰应用 \(\int_{\RR}\frac{dx}{x^2+a^2}=\frac{\pi}{a}\)（\(a>0\)）并作平移得 \(L^1\) 恒等式。
对 \(L^2\) 量，利用 \(H\ge0\) 从而
\(
\|H\|_2^2=\int_{\RR}H(x)^2dx\ge \sum_{j}\int_{\RR}(\text{第 }j\text{峰})^2dx
\)，
而
\(
\int_{\RR}\frac{dx}{(x^2+a^2)^2}=\frac{\pi}{2a^3}
\)
给出每峰平方积分为 \(8\pi m_j^2 \delta_j^2/(1+\delta_j)^3\)。
最后由 Hölder 不等式 \(\|H\|_2^2\le \|H\|_\infty\|H\|_1\) 即得峰值下界。证毕。
\end{proof}

\begin{corollary}[缺陷计数的双侧可审计界：由质量与峰值同时夹逼 \texorpdfstring{$\kappa$}{kappa}]\label{cor:xi-scan-defect-degree-bounded-by-mass-peak}
令 \(\delta_{\min}:=\min_j\delta_j\)、\(\delta_{\max}:=\max_j\delta_j\)。
则缺陷总重数满足两条可审计上界：
\[
\boxed{\;
\kappa\ \le\ \frac{1+\delta_{\max}}{4\pi\,\delta_{\min}}\ \|H\|_{L^1(\RR)},\qquad
\kappa\ \le\ \frac{(1+\delta_{\max})^{2}}{4\,\delta_{\min}}\ \|H\|_{L^\infty(\RR)}\;}
\]
\end{corollary}

\begin{proof}
由命题 \ref{prop:xi-scan-profile-lp-peak-witness} 的 \(L^1\) 恒等式并用
\(
\frac{\delta_j}{1+\delta_j}\ge\frac{\delta_{\min}}{1+\delta_{\max}}
\)
即得第一条。
第二条由点位估计：对任意 \(j\)，
\[
H(\gamma_j)\ge m_j\frac{4\delta_j}{(1+\delta_j)^2}\ge m_j\frac{4\delta_{\min}}{(1+\delta_{\max})^2},
\]
故 \(m_j\le \frac{(1+\delta_{\max})^2}{4\delta_{\min}}\|H\|_\infty\)，求和得结论。证毕。
\end{proof}

\begin{proposition}[Lorentzian 刚性：缺陷峰的幅度--半宽代数约束与可核验一致性检验]\label{prop:xi-lorentzian-rigidity-amplitude-halfwidth}
\leanverified{paper\_xi\_lorentzian\_rigidity\_amplitude\_halfwidth}
对单个点状缺陷原子 \((\gamma,\delta,m)\) 的共动剖面核
\[
H_{(\gamma,\delta,m)}(x)
:=m\,\frac{4\delta}{(x-\gamma)^2+(1+\delta)^2},
\]
其峰值 \(A:=H_{(\gamma,\delta,m)}(\gamma)\) 与半高半宽 \(w\)（定义为满足 \(H_{(\gamma,\delta,m)}(\gamma\pm w)=\tfrac12A\) 的正数）满足严格刚性关系
\[
\boxed{\;
w=1+\delta,\qquad
\frac{A}{m}=\frac{4(w-1)}{w^2}\in(0,1]\; }.
\]
因此在有限缺陷混合情形中，任何被制度判定为“孤立可分离峰”的局部结构必须满足上述幅度--半宽一致性；该一致性提供一个可离线复核的几何验伪准则。
并且由 \(a:=A/m\in(0,1]\) 可给出闭式反演
\[
\boxed{\;
w=\frac{2\bigl(1\mp\sqrt{1-a}\bigr)}{a},\qquad \delta=w-1.\;}
\]
在端点薄层/近临界口径（例如先验 \(\delta\le 1\)，等价 \(w\le 2\)）下取“\(-\)”分支即得到唯一恢复。
\end{proposition}

\begin{proof}
直接计算得
\(
A=m\frac{4\delta}{(1+\delta)^2}
\)。
令 \(H(\gamma\pm w)=\tfrac12A\) 并化简即得
\(
(\gamma\pm w-\gamma)^2+(1+\delta)^2=2(1+\delta)^2
\)，故 \(w=1+\delta\)。
再令 \(a=A/m\) 并以 \(w=1+\delta\) 代入可得二次方程
\(
a w^2-4w+4=0
\)，解之即得所示反演两分支；分支选择由先验口径（例如 \(\delta\le 1\)）固定。证毕。
\end{proof}

\begin{proposition}[基点扫描剖面的总质量恒等式：轴外偏离的全局可积不变量（48）]\label{prop:xi-scan-density-spectral-shift-phase-derivative}
\leanverified{paper\_xi\_scan\_density\_spectral\_shift\_phase\_derivative}
在有限点状缺陷口径下，定义上半平面内函数\footnote{据内部文稿 \cite{Internal2025ResolutionFoldingPhiPi}（2025）在“相位坐标门—有限窗审计—orbit--Levi 刚性”并置且不增设额外公理的口径下，该全质量恒等式与其 Fourier--Laplace 指纹、无穷远矩链一起构成可独立复核的扫描层析不变量链。}
\[
S_{\mathrm{HP}}(z)
:=\prod_{j=1}^{\kappa}\left(\frac{z-\gamma_j-\mathrm{i}(1+\delta_j)}{z-\gamma_j+\mathrm{i}(1+\delta_j)}\right)^{m_j},
\qquad \Im z>0.
\]
则其边界值满足 \(|S_{\mathrm{HP}}(x)|=1\)（a.e. \(x\in\RR\)），并具有相位导数公式
\[
\boxed{\;
\frac{d}{dx}\arg S_{\mathrm{HP}}(x)
=2\sum_{j=1}^{\kappa}m_j\,\frac{1+\delta_j}{(x-\gamma_j)^2+(1+\delta_j)^2}\; }.
\]
同时扫描剖面 \(H\)（\eqref{eq:xi-comoving-scan-profile-H}）满足严格加权谱移位表示
\[
\boxed{\;
H(x)
=\sum_{j=1}^{\kappa}m_j\,\frac{2\delta_j}{1+\delta_j}\,
\frac{d}{dx}\arg\!\left(\frac{x-\gamma_j-\mathrm{i}(1+\delta_j)}{x-\gamma_j+\mathrm{i}(1+\delta_j)}\right)\; }.
\]
特别地，其全质量为
\[
\boxed{\;\int_{\RR}H(x)\,\dd x
=4\pi\sum_{j=1}^{\kappa}m_j\,\frac{\delta_j}{1+\delta_j}.\;}
\]
\end{proposition}

\begin{proof}
对单因子 \(s_{j}(z):=\frac{z-\gamma_j-\mathrm{i}(1+\delta_j)}{z-\gamma_j+\mathrm{i}(1+\delta_j)}\)，在实轴上有 \(|s_j(x)|=1\) 且
\[
\frac{d}{dx}\log s_j(x)
=\frac{1}{x-\gamma_j-\mathrm{i}(1+\delta_j)}-\frac{1}{x-\gamma_j+\mathrm{i}(1+\delta_j)}
=\frac{2\mathrm{i}(1+\delta_j)}{(x-\gamma_j)^2+(1+\delta_j)^2},
\]
故 \(\frac{d}{dx}\arg s_j(x)=2\frac{1+\delta_j}{(x-\gamma_j)^2+(1+\delta_j)^2}\)；
乘上重数 \(m_j\) 并求和即得相位导数公式。
第二式由 \(H\) 的定义与上式逐项比对得到。
全质量恒等式由 \(\int_{\RR}\frac{dx}{(x-\gamma)^2+a^2}=\frac{\pi}{a}\) 直接推出。证毕。
\end{proof}

\paragraph{缺陷熵泛函：规范化面积、坐标不变性与指数--熵分解}
上式把扫描剖面写成相位流的加权密度，从而自然诱导一个仅依赖缺陷自由度、且对坐标规一化稳定的不变量：无量纲缺陷熵（或缺陷通量）。

\begin{definition}[无量纲缺陷熵/缺陷通量]\label{def:xi-defect-entropy-functional}
在有限点状缺陷口径下，定义缺陷熵（无量纲缺陷通量）
\begin{equation}\label{eq:xi-defect-entropy-functional}
\boxed{\;
S_{\mathrm{def}}
:=\frac{1}{4\pi}\int_{\RR}H(x)\,\dd x\;}
\end{equation}
其中 \(H\) 为 \eqref{eq:xi-comoving-scan-profile-H} 的共动扫描剖面。
\end{definition}

\begin{proposition}[缺陷熵的闭式表达与坐标不变性]\label{prop:xi-defect-entropy-closed-form-invariant}
\leanverified{paper\_xi\_defect\_entropy\_closed\_form\_invariant}
在上述设定下，
\begin{equation}\label{eq:xi-defect-entropy-closed-form}
\boxed{\;
S_{\mathrm{def}}
=\sum_{j=1}^{\kappa}m_j\,\frac{\delta_j}{1+\delta_j}\ \in\ (0,\kappa_{\mathrm{ind}}),\qquad
\kappa_{\mathrm{ind}}:=\sum_{j=1}^{\kappa}m_j.\;}
\end{equation}
并且对任意保持端点基准的 Möbius 重新规一化（例如改变 Cayley 标尺、从而一致重标定各处出现的共形平方常数；参见注 \ref{rem:dephys-cayley-scale-gauge-four}），\(H\) 的整体系数与坐标标尺发生共同重标定，而式 \eqref{eq:xi-defect-entropy-functional} 的无量纲组合 \(S_{\mathrm{def}}\) 保持不变。
\end{proposition}

\begin{proof}
由命题 \ref{prop:xi-scan-density-spectral-shift-phase-derivative} 的全质量恒等式
\(
\int_{\RR}H(x)\,\dd x
=4\pi\sum_{j}m_j\frac{\delta_j}{1+\delta_j}
\)
立得 \eqref{eq:xi-defect-entropy-closed-form}。
其中每一项满足 \(0<\frac{\delta_j}{1+\delta_j}<1\)，故 \(0<S_{\mathrm{def}}<\sum_j m_j=\kappa_{\mathrm{ind}}\)。
坐标不变性来自于：\(S_{\mathrm{def}}\) 由闭式表达 \(\sum_j m_j\frac{\delta_j}{1+\delta_j}\) 给出，而该比值正是端点 Poisson/Busemann 权重在规一化标尺下的无量纲组合；任何保持端点基准的共形重标定只会一致改变“单位选择”（对应 \(b_{-1}=\log B^{-1}\) 的加性零点），不会改变上述无量纲和。证毕。
\end{proof}

\begin{proposition}[指数--熵分解：整数缺陷指数与连续缺陷熵的严格分离]\label{prop:xi-index-entropy-decomposition}
\leanverified{paper\_xi\_index\_entropy\_decomposition}
令 \(S_{\mathrm{HP}}\) 为命题 \ref{prop:xi-scan-density-spectral-shift-phase-derivative} 的上半平面内函数。
记其边界相位流密度（谱移位密度）
\begin{equation}\label{eq:xi-spectral-shift-density-W}
\boxed{\;
W(x):=\frac{1}{2\pi}\frac{d}{dx}\arg S_{\mathrm{HP}}(x)\qquad(x\in\RR)\;}
\end{equation}
（取任一连续分支；改变分支不影响其导数）。
则其总质量给出整数缺陷指数
\begin{equation}\label{eq:xi-winding-index-kappa-ind}
\boxed{\;
\int_{\RR}W(x)\,\dd x=\kappa_{\mathrm{ind}}:=\sum_{j=1}^{\kappa}m_j\ \in\ \NN.\;}
\end{equation}
并且缺陷熵满足严格分解律
\begin{equation}\label{eq:xi-defect-entropy-index-decomposition}
\boxed{\;
S_{\mathrm{def}}
=\kappa_{\mathrm{ind}}-\sum_{j=1}^{\kappa}\frac{m_j}{1+\delta_j}\; }.
\end{equation}
因此制度允许的缺陷信息天然分裂为两类：\(\kappa_{\mathrm{ind}}\) 作为拓扑计数型整数指数（绕数不变量），以及 \(S_{\mathrm{def}}\) 作为随 \(\delta\) 连续变化的熵势（可见折损）。
\end{proposition}

\begin{proof}
由命题 \ref{prop:xi-scan-density-spectral-shift-phase-derivative}，有
\[
\frac{d}{dx}\arg S_{\mathrm{HP}}(x)
=\sum_{j=1}^{\kappa}m_j\frac{d}{dx}\arg s_j(x),
\qquad
s_j(x)=\frac{x-\gamma_j-\mathrm{i}(1+\delta_j)}{x-\gamma_j+\mathrm{i}(1+\delta_j)}.
\]
而 \(\arg s_j(x)\) 的相位变差在 \(x\in(-\infty,\infty)\) 上恰为 \(2\pi\)（按取向计；取绝对值则为 \(2\pi\)），故
\(
\frac{1}{2\pi}\int_{\RR}\frac{d}{dx}\arg s_j(x)\,\dd x=1
\)；
乘上重数并求和即得 \eqref{eq:xi-winding-index-kappa-ind}。
\par
式 \eqref{eq:xi-defect-entropy-index-decomposition} 则由
\(
\frac{\delta}{1+\delta}=1-\frac{1}{1+\delta}
\)
逐项代入 \eqref{eq:xi-defect-entropy-closed-form} 得到。证毕。
\end{proof}

\begin{corollary}[规范化“面积律”与熵塌缩/饱和二相图]\label{cor:xi-defect-entropy-area-law-phase-diagram}
在上述设定下，
\[
\boxed{\;
S_{\mathrm{def}}
=\int_{\RR}\frac{1}{4\pi}H(x)\,\dd x
=\int_{\RR}\sum_{j=1}^{\kappa}\frac{\delta_j}{1+\delta_j}\,m_j\,
\frac{1}{2\pi}\frac{d}{dx}\arg s_j(x)\,\dd x\;}
\]
从而 \(S_{\mathrm{def}}\) 可视作一条“加权谱移位密度”的规范化面积。
此外，对单原子缺陷贡献为 \(m\frac{\delta}{1+\delta}\)，故端点薄层与远离薄层的两端极限给出严格二相图：
\[
\boxed{\;
\delta\downarrow 0\ \Longrightarrow\ S_{\mathrm{def}}\downarrow 0,\qquad
\delta\uparrow\infty\ \Longrightarrow\ S_{\mathrm{def}}\uparrow \kappa_{\mathrm{ind}}.\;}
\]
\end{corollary}

\paragraph{缺陷熵的双曲几何律：面积、horocycle 长度与尺度指纹}
本段给出若干仅依赖 \(\{(\gamma_j,\delta_j,m_j)\}\) 的标量后果：它们把
\(\displaystyle S_{\mathrm{def}}=\sum_j m_j\frac{\delta_j}{1+\delta_j}\)
严格识别为双曲面积与 horocycle 长度亏损，并给出由线性通量与整数指数所强制的最优可行区间；
此外，通过尺度粗粒化得到一条完全单调的 Stieltjes 轨道，从而将深度谱进一步压缩为可审计的单变量指纹。

\begin{proposition}[缺陷熵的双曲面积律与 \texorpdfstring{$4\pi$}{4pi} 规一化]\label{prop:xi-defect-entropy-hyperbolic-area-law-4pi}
在上半平面模型 \(\mathbb{H}=\{x+\mathrm{i}y:\ y>0\}\) 取双曲度量
\(
ds^2=(dx^2+dy^2)/y^2
\)，面积元
\(
\dd A_{\mathbb{H}}=\frac{\dd x\,\dd y}{y^2}
\)。
对每个点状缺陷参数 \((\gamma_j,\delta_j,m_j)\)（\(\delta_j>0\)）定义单位宽竖条域
\[
R_j:=
\bigl[\gamma_j-\tfrac12,\gamma_j+\tfrac12\bigr]\times\bigl[1,1+\delta_j\bigr]\subset\mathbb{H}.
\]
则其双曲面积满足严格恒等式
\begin{equation}\label{eq:xi-hyperbolic-area-vertical-strip}
\boxed{\;
A_{\mathbb{H}}(R_j)
:=\int_{R_j}\dd A_{\mathbb{H}}
=\frac{\delta_j}{1+\delta_j}\; }.
\end{equation}
从而缺陷熵具有双曲面积律
\begin{equation}\label{eq:xi-defect-entropy-hyperbolic-area-law}
\boxed{\;
S_{\mathrm{def}}
=\sum_{j=1}^{\kappa}m_j\,A_{\mathbb{H}}(R_j)\; }.
\end{equation}
另一方面，共动扫描剖面 \(H\)（\eqref{eq:xi-comoving-scan-profile-H}）满足逐原子恒等式
\begin{equation}\label{eq:xi-atom-mass-4pi-area-alignment}
\boxed{\;
\int_{\RR} m_j\,\frac{4\delta_j}{(x-\gamma_j)^2+(1+\delta_j)^2}\,\dd x
=4\pi\,m_j\,A_{\mathbb{H}}(R_j)\; }.
\end{equation}
因此
\(
\frac{1}{4\pi}\int_{\RR}H(x)\dd x
=S_{\mathrm{def}}
\)
与定义 \ref{def:xi-defect-entropy-functional} 一致；
“\(4\pi\)”的作用即在于把边界 Poisson 迹的欧氏积分与双曲面积在同一规一化下对齐。
\leanverified{paper\_xi\_defect\_entropy\_hyperbolic\_area\_law\_4pi}
\end{proposition}

\begin{proof}
由 \(\dd A_{\mathbb{H}}=\frac{\dd x\,\dd y}{y^2}\) 直接计算
\[
A_{\mathbb{H}}(R_j)
=\int_{\gamma_j-1/2}^{\gamma_j+1/2}\int_{1}^{1+\delta_j}y^{-2}\,\dd y\,\dd x
=\int_{1}^{1+\delta_j}y^{-2}\,\dd y
=1-\frac{1}{1+\delta_j}
=\frac{\delta_j}{1+\delta_j},
\]
得到 \eqref{eq:xi-hyperbolic-area-vertical-strip}；代入 \eqref{eq:xi-defect-entropy-closed-form} 即得 \eqref{eq:xi-defect-entropy-hyperbolic-area-law}。
\par
对 \eqref{eq:xi-atom-mass-4pi-area-alignment}，用标准积分
\(
\int_{\RR}\frac{\dd x}{(x-\gamma)^2+a^2}=\frac{\pi}{a}
\)（\(a>0\)）并取 \(a=1+\delta_j\) 即得
\[
\int_{\RR} m_j\,\frac{4\delta_j}{(x-\gamma_j)^2+(1+\delta_j)^2}\,\dd x
=4\pi\,m_j\,\frac{\delta_j}{1+\delta_j}
=4\pi\,m_j\,A_{\mathbb{H}}(R_j).
\]
证毕。
\end{proof}

\begin{corollary}[边界长度亏损公式：熵即 horocycle 长度的净耗散]\label{cor:xi-defect-entropy-horocycle-length-deficit}
\leanverified{paper\_xi\_defect\_entropy\_horocycle\_length\_deficit}
在同一双曲度量下，设高度 \(y\) 处欧氏长度为 \(1\) 的水平线段之双曲长度为
\[
L(y):=\int_{-1/2}^{1/2}\frac{\dd x}{y}=\frac{1}{y}.
\]
则对每个条域 \(R_j\)（命题 \ref{prop:xi-defect-entropy-hyperbolic-area-law-4pi}）有
\[
A_{\mathbb{H}}(R_j)=L(1)-L(1+\delta_j)=1-\frac{1}{1+\delta_j}.
\]
因此缺陷熵满足严格恒等式
\begin{equation}\label{eq:xi-defect-entropy-horocycle-length-deficit}
\boxed{\;
S_{\mathrm{def}}
=\sum_{j=1}^{\kappa}m_j\bigl(L(1)-L(1+\delta_j)\bigr)
=\kappa_{\mathrm{ind}}-\sum_{j=1}^{\kappa}\frac{m_j}{1+\delta_j}\; }.
\end{equation}
该恒等式把 \(S_{\mathrm{def}}\) 解释为：在底部 horocycle（\(y=1\)）上，每个缺陷将同宽条域提升到高度 \(1+\delta_j\) 所产生的双曲长度亏损；熵即亏损的可加总量。
\end{corollary}

\begin{proof}
由 \(A_{\mathbb{H}}(R_j)=\int_{1}^{1+\delta_j}y^{-2}\dd y\) 与 \(L(y)=1/y\) 的定义，得
\(
A_{\mathbb{H}}(R_j)=L(1)-L(1+\delta_j)
\)。
将该恒等式代入
\(
S_{\mathrm{def}}=\sum_j m_jA_{\mathbb{H}}(R_j)
\)
并用 \(L(1)=1\) 即得 \eqref{eq:xi-defect-entropy-horocycle-length-deficit}；
其与 \eqref{eq:xi-defect-entropy-index-decomposition} 一致。证毕。
\end{proof}

\begin{proposition}[熵--通量--指数的最优可行区间与调和平均下界]\label{prop:xi-defect-entropy-feasible-interval-flux-index}
\leanverified{paper\_xi\_defect\_entropy\_feasible\_interval\_flux\_index}
令线性通量（深度总量）
\[
\Phi:=\sum_{j=1}^{\kappa}m_j\,\delta_j
\qquad\bigl(=\ M_0^{\mathrm{def}}\ \text{于命题 \ref{prop:xi-poisson-coarsegraining-nohair} 的记号}\bigr),
\]
并记整数缺陷指数 \(\kappa_{\mathrm{ind}}:=\sum_{j=1}^{\kappa}m_j\)。
则缺陷熵满足锐界
\begin{equation}\label{eq:xi-defect-entropy-feasible-interval}
\boxed{\;
\frac{\Phi}{1+\Phi}
\ \le\
S_{\mathrm{def}}
\ \le\
\frac{\kappa_{\mathrm{ind}}\Phi}{\kappa_{\mathrm{ind}}+\Phi}\; }.
\end{equation}
上界在“均匀深度”（按重数展开后 \(\delta\) 常值）情形达到；下界在“极端集中”情形达到其下确界（把全部深度集中到一个单位重数原子，其余深度趋于 \(0\)）。
特别地，当 \(\Phi>0\) 时，上界等价改写为调和平均型不等式
\begin{equation}\label{eq:xi-defect-entropy-harmonic-mean-bound}
\boxed{\;
\frac{1}{S_{\mathrm{def}}}\ \ge\ \frac{1}{\Phi}+\frac{1}{\kappa_{\mathrm{ind}}}\; }.
\end{equation}
\end{proposition}

\begin{proof}
令 \(f(\delta):=\frac{\delta}{1+\delta}\)，则 \(f\) 在 \((0,\infty)\) 上递增且严格凹：
\(
f''(\delta)=-2(1+\delta)^{-3}<0
\)。
\par
\textbf{上界：}
以权重 \(w_j:=m_j/\kappa_{\mathrm{ind}}\)（\(\sum_j w_j=1\)）作 Jensen 不等式，得
\[
\frac{S_{\mathrm{def}}}{\kappa_{\mathrm{ind}}}
=\sum_{j}w_j f(\delta_j)
\le f\!\Bigl(\sum_j w_j\delta_j\Bigr)
=f\!\Bigl(\frac{\Phi}{\kappa_{\mathrm{ind}}}\Bigr)
=\frac{\Phi}{\kappa_{\mathrm{ind}}+\Phi},
\]
从而 \(S_{\mathrm{def}}\le \frac{\kappa_{\mathrm{ind}}\Phi}{\kappa_{\mathrm{ind}}+\Phi}\)。
\par
\textbf{下界：}
由于 \(m_j\in\NN\)（本节设定），有 \(\delta_j\le m_j\delta_j\le \Phi\)，故
\(
\frac{\delta_j}{1+\delta_j}\ge \frac{\delta_j}{1+\Phi}
\)。
逐项乘以 \(m_j\) 并求和得
\[
S_{\mathrm{def}}
=\sum_j m_j\frac{\delta_j}{1+\delta_j}
\ge \sum_j m_j\frac{\delta_j}{1+\Phi}
=\frac{\Phi}{1+\Phi}.
\]
最后，\eqref{eq:xi-defect-entropy-harmonic-mean-bound} 由上界取倒数并整理得到。证毕。
\end{proof}

\begin{proposition}[Schur--凹性：缺陷熵对深度向量的主序最优化原则]\label{prop:xi-defect-entropy-schur-concavity}
\leanverified{paper\_xi\_defect\_entropy\_schur\_concavity}
固定 \(\kappa_{\mathrm{ind}}\ge 2\) 并把深度按重数展开为向量
\(\boldsymbol{\delta}=(\delta_1,\dots,\delta_{\kappa_{\mathrm{ind}}})\in(0,\infty)^{\kappa_{\mathrm{ind}}}\)。
定义
\[
\mathcal{S}(\boldsymbol{\delta})
:=\sum_{i=1}^{\kappa_{\mathrm{ind}}}\frac{\delta_i}{1+\delta_i}.
\]
则 \(\mathcal{S}\) 为严格 Schur--凹函数：若 \(\boldsymbol{\delta}\) 在主序意义上\textbf{更集中}（即 \(\boldsymbol{\delta}\succ \boldsymbol{\delta}'\) 且二者不同于排列），则
\[
\mathcal{S}(\boldsymbol{\delta})\ <\ \mathcal{S}(\boldsymbol{\delta}').
\]
因此在任意线性约束 \(\sum_{i=1}^{\kappa_{\mathrm{ind}}}\delta_i=\Phi\) 下，
\(\mathcal{S}\) 的最大值由等深度配置 \(\delta_i\equiv \Phi/\kappa_{\mathrm{ind}}\) 达到，而其下确界由极端集中配置（一个分量趋于 \(\Phi\) 其余分量趋于 \(0\)）逼近。
该主序原则与命题 \ref{prop:xi-defect-entropy-feasible-interval-flux-index} 的“上界在均匀深度取等号、下界在极端集中取下确界”完全一致，并将其提升为结构性最优化律。
\end{proposition}
\begin{proof}
函数 \(f(\delta)=\delta/(1+\delta)\) 在 \((0,\infty)\) 上严格凹（见命题 \ref{prop:xi-defect-entropy-feasible-interval-flux-index} 的计算），
故由 Karamata 不等式（严格情形）知 \(\sum_i f(\delta_i)\) 为严格 Schur--凹。
在固定总和约束下，等深度配置为主序最小元，从而给出最大值；极端集中为主序最大端点，从而给出下确界。证毕。
\end{proof}

% Peak-compressibility extremal law under fixed defect entropy.

\begin{theorem}\label{thm:xi-defect-entropy-peak-compressibility-extremal}
\leanverified{paper\_xi\_defect\_entropy\_peak\_compressibility\_extremal}
沿用共动扫描剖面 \(H\) 的点状缺陷参数化：给定整数重数 \(m_j\in\ZZ_{>0}\)、深度参数 \(\delta_j>0\) 与中心 \(\gamma_j\in\RR\)，令
\[
H(x)=\sum_{j=1}^{J} m_j\,\frac{4\delta_j}{(x-\gamma_j)^2+(1+\delta_j)^2},
\qquad x\in\RR,
\]
并记缺陷指数
\(
\kappa_{\mathrm{ind}}:=\sum_{j=1}^{J}m_j
\)
与缺陷熵
\[
S_{\mathrm{def}}
:=\frac{1}{4\pi}\int_{\RR}H(x)\,\dd x
=\sum_{j=1}^{J}m_j\frac{\delta_j}{1+\delta_j}\in(0,\kappa_{\mathrm{ind}}).
\]
在固定 \((\kappa_{\mathrm{ind}},S_{\mathrm{def}})\) 下，考虑峰值下确界
\[
\mathcal P(\kappa_{\mathrm{ind}},S_{\mathrm{def}})
:=\inf \|H\|_{L^\infty(\RR)},
\]
其中下确界在所有满足上述约束的参数族 \(\{(m_j,\delta_j,\gamma_j)\}\) 上取。
设
\[
S_{\mathrm{def}}=M+r,\qquad M:=\lfloor S_{\mathrm{def}}\rfloor,\quad r\in[0,1),
\]
并定义
\[
a_\ast:=\min\Bigl\{\frac{r}{\kappa_{\mathrm{ind}}-M},\ \frac{1-r}{M+1}\Bigr\}\in\Bigl[0,\frac12\Bigr].
\]
则有显式极值公式
\[
\mathcal P(\kappa_{\mathrm{ind}},S_{\mathrm{def}})=4a_\ast(1-a_\ast).
\]
此外，上述下确界可由两类二值化构型在分离极限 \(|\gamma_j-\gamma_{j'}|\to\infty\) 中逼近达到：按重数展开后的 \(\kappa_{\mathrm{ind}}\) 个
\(
g:=\delta/(1+\delta)\in(0,1)
\)
仅取两端区间 \([0,a_\ast]\cup[1-a_\ast,1]\) 的端点型取值，并使 \(\sum g=S_{\mathrm{def}}\)。
\end{theorem}
\begin{proof}
记 \(g_j:=\delta_j/(1+\delta_j)\in(0,1)\)，则
\(
\frac{4\delta_j}{(1+\delta_j)^2}=4g_j(1-g_j)
\)。
对每个中心点 \(x=\gamma_j\)，有
\[
H(\gamma_j)\ge m_j\frac{4\delta_j}{(1+\delta_j)^2}=4m_jg_j(1-g_j),
\]
从而 \(\|H\|_\infty\ge 4\max_\ell g_\ell(1-g_\ell)\)，其中 \((g_\ell)_{\ell=1}^{\kappa_{\mathrm{ind}}}\) 表示将 \(g_j\) 按重数 \(m_j\) 展开后的序列。
于是
\[
\mathcal P(\kappa_{\mathrm{ind}},S_{\mathrm{def}})
\ge
4\cdot \inf\Bigl\{\max_{1\le \ell\le \kappa_{\mathrm{ind}}} g_\ell(1-g_\ell):\ g_\ell\in(0,1),\ \sum_{\ell=1}^{\kappa_{\mathrm{ind}}}g_\ell=S_{\mathrm{def}}\Bigr\}.
\]
设 \(\max_\ell g_\ell(1-g_\ell)\le M_0\le 1/4\)，令 \(a\in[0,1/2]\) 为 \(a(1-a)=M_0\) 的小根。
则对任意 \(g\in(0,1)\)，
\(
g(1-g)\le M_0\iff g\in[0,a]\cup[1-a,1]
\)。
若其中恰有 \(m\) 个落在 \([1-a,1]\)，其余 \(\kappa_{\mathrm{ind}}-m\) 个落在 \([0,a]\)，则总和满足
\[
m(1-a)\le \sum_{\ell=1}^{\kappa_{\mathrm{ind}}}g_\ell \le m+(\kappa_{\mathrm{ind}}-m)a.
\]
令 \(M=\lfloor S_{\mathrm{def}}\rfloor\) 与 \(S_{\mathrm{def}}=M+r\)。
可行性要求 \(S_{\mathrm{def}}\) 落入上述区间并对某个整数 \(m\) 成立。
随着 \(a\) 增大，区间端点关于 \(m\) 的最小可行候选只可能发生在 \(m=M\) 或 \(m=M+1\)；
分别解端点等式得到
\[
a\ge \frac{r}{\kappa_{\mathrm{ind}}-M},
\qquad
a\ge \frac{1-r}{M+1}.
\]
因此最小可行 \(a\) 等于 \(a_\ast\)，从而最小可能 \(M_0\) 为 \(a_\ast(1-a_\ast)\)，给出下界
\(
\mathcal P\ge 4a_\ast(1-a_\ast)
\)。
\par
对上界构造，取展开序列中 \(M\) 个 \(g\) 逼近 \(1\)，其余 \(\kappa_{\mathrm{ind}}-M\) 个取定为 \(a_\ast\)（或在 \(m=M+1\) 的对偶构型中取 \(M+1\) 个 \(g\equiv 1-a_\ast\) 且其余逼近 \(0\)），并令对应中心 \(\gamma\) 两两远离。
由于每个核为 \(O(|x-\gamma|^{-2})\) 衰减，分离极限下不同核在彼此中心处的尾部叠加可任意压小，故 \(\|H\|_\infty\) 可逼近 \(4a_\ast(1-a_\ast)\)。
结合下界即得等式。证毕。
\end{proof}

\begin{corollary}\label{cor:xi-defect-entropy-mass-peak-decoupling}
\leanverified{paper\_xi\_defect\_entropy\_mass\_peak\_decoupling}
在定理 \ref{thm:xi-defect-entropy-peak-compressibility-extremal} 的设定下，对任意 \(0<S_{\mathrm{def}}<\kappa_{\mathrm{ind}}\) 总有
\[
\mathcal P(\kappa_{\mathrm{ind}},S_{\mathrm{def}})\le \frac{8}{\kappa_{\mathrm{ind}}+1}.
\]
因此当 \(\kappa_{\mathrm{ind}}\to\infty\) 且 \(S_{\mathrm{def}}\) 固定（或至多作 \(O(\kappa_{\mathrm{ind}})\) 增长）时，\(\mathcal P(\kappa_{\mathrm{ind}},S_{\mathrm{def}})=O(\kappa_{\mathrm{ind}}^{-1})\)。
\end{corollary}
\begin{proof}
由定理 \ref{thm:xi-defect-entropy-peak-compressibility-extremal}，
\(
\mathcal P=4a_\ast(1-a_\ast)\le 4a_\ast
\)。
又 \(r\in[0,1)\) 给出
\[
a_\ast=\min\Bigl\{\frac{r}{\kappa_{\mathrm{ind}}-M},\frac{1-r}{M+1}\Bigr\}
\le \min\Bigl\{\frac{1}{\kappa_{\mathrm{ind}}-M},\frac{1}{M+1}\Bigr\}
\le \frac{2}{\kappa_{\mathrm{ind}}+1},
\]
从而 \(\mathcal P\le 8/(\kappa_{\mathrm{ind}}+1)\)。证毕。
\end{proof}

\endinput


\begin{proposition}[统一伸缩的熵重整化不等式：\texorpdfstring{$\delta\mapsto m\delta$}{delta -> m delta} 的 Jensen 包络]\label{prop:xi-defect-entropy-uniform-scaling-jensen-envelope}
\leanverified{paper\_xi\_defect\_entropy\_uniform\_scaling\_jensen\_envelope}
令
\(
p_j:=\delta_j/(1+\delta_j)\in(0,1)
\)，并记
\(
S:=S_{\mathrm{def}}=\sum_j m_j p_j
\)、
\(
\kappa:=\kappa_{\mathrm{ind}}=\sum_j m_j
\)。
对任意 \(m>0\) 施加统一伸缩 \(\delta_j\mapsto m\delta_j\)，并定义新熵
\[
S_m:=\sum_j m_j\,\frac{m\delta_j}{1+m\delta_j}.
\]
则在 \(p\)-坐标下该变换由 Möbius 变换
\begin{equation}\label{eq:xi-defect-entropy-scaling-mobius}
\boxed{\;
\Psi_m(p):=\frac{mp}{1+(m-1)p}\qquad(p\in[0,1])\;}
\end{equation}
给出：\(S_m=\sum_j m_j\,\Psi_m(p_j)\)。
并且：
\begin{enumerate}
  \item 若 \(m>1\)，则 \(\Psi_m\) 在 \([0,1]\) 上严格凹，故由 Jensen 不等式得
  \[
  \boxed{\;
  S_m\ \le\ \kappa\,\Psi_m\!\left(\frac{S}{\kappa}\right),\;}
  \]
  且等号当且仅当所有 \(p_j\)（按重数展开）相等；
  \item 若 \(0<m<1\)，则 \(\Psi_m\) 在 \([0,1]\) 上严格凸，故同理得到
  \[
  \boxed{\;
  S_m\ \ge\ \kappa\,\Psi_m\!\left(\frac{S}{\kappa}\right),\;}
  \]
  且等号当且仅当所有 \(p_j\)（按重数展开）相等。
\end{enumerate}
因此，“统一尺度变化下熵的可达域”被压缩为单变量 Möbius 包络，而“均匀深度”是唯一饱和构型。
\end{proposition}

\begin{proof}
由 \(p=\delta/(1+\delta)\) 解出 \(\delta=p/(1-p)\)，代入 \(m\delta\) 并化简得
\(\frac{m\delta}{1+m\delta}=\Psi_m(p)\)，从而 \(S_m=\sum_j m_j\Psi_m(p_j)\)。
计算
\[
\Psi_m'(p)=\frac{m}{(1+(m-1)p)^2},
\qquad
\Psi_m''(p)=-\frac{2m(m-1)}{(1+(m-1)p)^3},
\]
故 \(m>1\) 时 \(\Psi_m''<0\)（严格凹），\(0<m<1\) 时 \(\Psi_m''>0\)（严格凸）。
再以权重 \(m_j/\kappa\) 应用 Jensen 即得两条不等式；严格取等条件由严格凸/凹的等号条件给出。证毕。
\end{proof}

\begin{proposition}[连续重整化流的方差缺口恒等式：偏离均匀性即回收强度]\label{prop:xi-defect-entropy-scaling-variance-gap-identity}
\leanverified{paper\_xi\_defect\_entropy\_scaling\_variance\_gap\_identity}
以对数尺度 \(s=\log m\) 参数化统一伸缩 \(\delta\mapsto e^{s}\delta\)，并令
\(
p_j(s):=\Psi_{e^{s}}(p_j(0))
\)、
\(
S(s):=\sum_j m_j p_j(s)
\)、
\(
\bar p(s):=S(s)/\kappa
\)。
则对所有 \(s\in\RR\) 有导数闭式
\begin{equation}\label{eq:xi-defect-entropy-logistic-variance-gap}
\boxed{\;
\frac{\dd}{\dd s}S(s)
=\sum_j m_j\,p_j(s)\bigl(1-p_j(s)\bigr)
=S(s)\Bigl(1-\frac{S(s)}{\kappa}\Bigr)-\kappa\,\mathrm{Var}_m\!\bigl(p(s)\bigr),\;}
\end{equation}
其中加权方差定义为
\[
\mathrm{Var}_m(p):=\frac{1}{\kappa}\sum_j m_j\bigl(p_j-\bar p\bigr)^2\ \ge\ 0.
\]
因此
\[
\boxed{\;
\frac{\dd}{\dd s}S(s)\ \le\ S(s)\Bigl(1-\frac{S(s)}{\kappa}\Bigr),\;}
\]
且严格不等号当且仅当深度非均匀（\(\mathrm{Var}_m(p(s))>0\)）。
该恒等式把“回收器强度”量化为一个可计算的缺口项：均匀化是唯一达到 logistic 速率上界的情形，而任何非均匀性都以方差在制度层面被额外抹平。
\end{proposition}

\begin{proof}
由 \eqref{eq:xi-defect-entropy-scaling-mobius} 的闭式可直接验证
\[
\frac{\partial}{\partial s}\Psi_{e^s}(p)
=\frac{e^sp}{(1+(e^s-1)p)^2}
=\Psi_{e^s}(p)\bigl(1-\Psi_{e^s}(p)\bigr),
\]
故 \(p_j'(s)=p_j(s)(1-p_j(s))\)，从而 \(\frac{\dd}{\dd s}S(s)=\sum_j m_j p_j(s)(1-p_j(s))\)。
再用恒等式
\[
\sum_j m_j p_j(1-p_j)
\;=\;
\kappa\,\bar p(1-\bar p)\;-\;\kappa\,\mathrm{Var}_m(p),
\]
并代入 \(\bar p=S/\kappa\)，即得 \eqref{eq:xi-defect-entropy-logistic-variance-gap}。证毕。
\end{proof}

\begin{corollary}[证书复杂度的双曲指数律：有限阶矩证书只覆盖对数半径]\label{cor:xi-certificate-complexity-hyperbolic-exponential-law}
设 \(a\in\mathbb{D}\) 为盘内点状缺陷（等价视界对数导数在 \(a\) 处出现极点；定理 \ref{thm:xi-disk-zero-pole-index-reduction}），并记其欧氏边界深度
\(
h:=1-|a|\in(0,1)
\)
以及从圆盘中心到 \(a\) 的双曲距离
\[
d:=\log\frac{1+|a|}{1-|a|}.
\]
若某一有限阶 Toeplitz--PSD 证书（或等价有限 CMV/Schur 证书）能够以制度意义下的稳定负方向见证该缺陷，则其阶数 \(N\) 必满足
\[
\boxed{\;
N\ \gtrsim\ h^{-1}\;}
\]
（见定理 \ref{thm:xi-toeplitz-length-vs-boundary-depth} 的粗尺度必要下界）。
由恒等式 \(h=\frac{2}{e^{d}+1}\) 得到等价的双曲指数下界
\[
\boxed{\;
N\ \gtrsim\ e^{d}\;}
\]
从而给定证书阶数 \(N\)，其可稳定覆盖的双曲距离至多为 \(d\lesssim \log N\) 的量级。
因此，有限阶矩证书的覆盖半径在双曲几何中仅按 \(\log N\) 增长，而证书资源随双曲深度呈指数膨胀。
\end{corollary}

\begin{proof}
由 \(d=\log\frac{1+|a|}{1-|a|}\) 解出
\(
|a|=\frac{e^{d}-1}{e^{d}+1}
\)，故
\(
h=1-|a|=\frac{2}{e^{d}+1}
\)。
将定理 \ref{thm:xi-toeplitz-length-vs-boundary-depth} 的粗尺度必要下界 \(N\gtrsim h^{-1}\) 代入即得 \(N\gtrsim e^{d}\)；
反向结论由 \(d=\log\frac{2-h}{h}\le \log\frac{2}{h}\) 与 \(h\gtrsim N^{-1}\) 直接推出。证毕。
\end{proof}

\begin{proposition}[双曲高度线性化：熵与顶边长度的指数互补与半群律]\label{prop:xi-defect-entropy-hyperbolic-height-semigroup}
\leanverified{paper\_xi\_defect\_entropy\_hyperbolic\_height\_semigroup}
引入双曲高度变量
\[
s_j:=\log(1+\delta_j)\qquad(1+\delta_j=e^{s_j}).
\]
则每个原子的熵贡献满足
\[
\frac{\delta_j}{1+\delta_j}=1-e^{-s_j},
\]
从而缺陷熵具有指数线性化表达
\begin{equation}\label{eq:xi-defect-entropy-hyperbolic-height-linearization}
\boxed{\;
S_{\mathrm{def}}=\sum_{j=1}^{\kappa}m_j\bigl(1-e^{-s_j}\bigr)\;}
\end{equation}
并且“顶边总长度”（亦即高度 \(y=1+\delta_j\) 处单位宽 horocycle 段的加权总长）
\[
L_{\mathrm{top}}
:=\sum_{j=1}^{\kappa}\frac{m_j}{1+\delta_j}
=\sum_{j=1}^{\kappa}m_j e^{-s_j}
=\kappa_{\mathrm{ind}}-S_{\mathrm{def}}.
\]
特别地，若对所有原子施加统一竖直平移 \(s_j\mapsto s_j+\tau\)（\(\tau\ge 0\)），则
\[
L_{\mathrm{top}}\mapsto e^{-\tau}L_{\mathrm{top}},
\qquad
S_{\mathrm{def}}\mapsto \kappa_{\mathrm{ind}}-e^{-\tau}\bigl(\kappa_{\mathrm{ind}}-S_{\mathrm{def}}\bigr),
\]
从而 \(\tau\) 生成一参数半群，并给出严格微分式
\[
\dd S_{\mathrm{def}}=\sum_{j=1}^{\kappa} m_j e^{-s_j}\,\dd s_j.
\]
\end{proposition}

\begin{proof}
由 \(1+\delta_j=e^{s_j}\) 得
\(
\frac{\delta_j}{1+\delta_j}=1-\frac{1}{1+\delta_j}=1-e^{-s_j}
\)，代入 \eqref{eq:xi-defect-entropy-closed-form} 即得 \eqref{eq:xi-defect-entropy-hyperbolic-height-linearization}。
其余各式由代数恒等
\(
L_{\mathrm{top}}=\sum_j m_j e^{-s_j}
\)
与
\(
S_{\mathrm{def}}=\kappa_{\mathrm{ind}}-L_{\mathrm{top}}
\)
逐项推出；微分式由对 \(s_j\) 逐项求导得到。证毕。
\end{proof}

\begin{corollary}[双曲高度预算的指数律：固定 \texorpdfstring{$\sum m_j\log(1+\delta_j)$}{sum m log(1+delta)} 时缺陷熵的最优上界]\label{cor:xi-defect-entropy-height-budget-exponential-law}
在命题 \ref{prop:xi-defect-entropy-hyperbolic-height-semigroup} 的设定下，令
\[
S_{\mathrm{ht}}:=\sum_{j=1}^{\kappa}m_j\log(1+\delta_j)=\sum_{j=1}^{\kappa}m_j s_j,
\qquad
\kappa_{\mathrm{ind}}:=\sum_{j=1}^{\kappa}m_j.
\]
则缺陷熵在固定高度预算 \(S_{\mathrm{ht}}\) 下满足锐上界
\begin{equation}\label{eq:xi-defect-entropy-height-budget-exp-law}
\boxed{\;
S_{\mathrm{def}}
\ \le\
\kappa_{\mathrm{ind}}\Bigl(1-e^{-S_{\mathrm{ht}}/\kappa_{\mathrm{ind}}}\Bigr).\;}
\end{equation}
并且当且仅当所有 \(s_j\)（按重数计）相等（等价于所有 \(\delta_j\) 相等）时取等号。
该不等式给出“面积/高度预算”约束下缺陷熵的指数型饱和律，并将最优态唯一化为均匀深度配置。
\end{corollary}
\begin{proof}
由 \eqref{eq:xi-defect-entropy-hyperbolic-height-linearization}，
\(
S_{\mathrm{def}}=\sum_j m_j g(s_j)
\)
其中 \(g(s)=1-e^{-s}\) 在 \([0,\infty)\) 上严格凹。
以权重 \(w_j:=m_j/\kappa_{\mathrm{ind}}\)（\(\sum_j w_j=1\)）作 Jensen 不等式，得
\[
\frac{S_{\mathrm{def}}}{\kappa_{\mathrm{ind}}}
=\sum_j w_j g(s_j)
\le g\!\Bigl(\sum_j w_j s_j\Bigr)
=1-\exp\!\Bigl(-\frac{S_{\mathrm{ht}}}{\kappa_{\mathrm{ind}}}\Bigr),
\]
即得 \eqref{eq:xi-defect-entropy-height-budget-exp-law}。
严格凹性给出等号条件：当且仅当 \(s_j\)（按重数计）全相等。证毕。
\end{proof}

\begin{theorem}[高度预算下缺陷熵的二阶精化：方差缺口项]\label{thm:xi-defect-entropy-height-budget-variance-gap}
\leanverified{paper\_xi\_defect\_entropy\_height\_budget\_variance\_gap}
在推论 \ref{cor:xi-defect-entropy-height-budget-exponential-law} 的设定下，令
\[
\bar s:=\frac{S_{\mathrm{ht}}}{\kappa_{\mathrm{ind}}},
\qquad
\mathrm{Var}_m(s):=\frac{1}{\kappa_{\mathrm{ind}}}\sum_{j=1}^{\kappa}m_j(s_j-\bar s)^2.
\]
则缺陷熵满足二阶精化上界
\[
\boxed{\;
S_{\mathrm{def}}
\ \le\
\kappa_{\mathrm{ind}}\Bigl(1-e^{-\bar s}\Bigr)
\ -\ \frac{\kappa_{\mathrm{ind}}}{2}\,e^{-\bar s}\,\mathrm{Var}_m(s)\; }.
\]
并且等号成立当且仅当 \(s_j\)（按重数计）全相等（等价于 \(\mathrm{Var}_m(s)=0\)）。
\end{theorem}
\begin{proof}
由命题 \ref{prop:xi-defect-entropy-hyperbolic-height-semigroup} 的线性化恒等式，
\[
\kappa_{\mathrm{ind}}-S_{\mathrm{def}}=\sum_{j=1}^{\kappa}m_j e^{-s_j}
=\kappa_{\mathrm{ind}}\,e^{-\bar s}\cdot \sum_{j=1}^{\kappa}\frac{m_j}{\kappa_{\mathrm{ind}}}\,e^{-(s_j-\bar s)}.
\]
令 \(x_j:=s_j-\bar s\)，则 \(\sum_j m_j x_j=0\)。
并且对任意实数 \(x\) 有基本不等式 \(e^{-x}\ge 1-x+\frac{x^2}{2}\)。
以权重 \(m_j/\kappa_{\mathrm{ind}}\) 加权求和并利用 \(\sum_j m_j x_j=0\)，得
\[
\sum_{j=1}^{\kappa}\frac{m_j}{\kappa_{\mathrm{ind}}}\,e^{-x_j}
\ge
1+\frac{1}{2}\sum_{j=1}^{\kappa}\frac{m_j}{\kappa_{\mathrm{ind}}}\,x_j^2
=1+\frac{1}{2}\mathrm{Var}_m(s).
\]
代回即得
\(
\kappa_{\mathrm{ind}}-S_{\mathrm{def}}
\ge
\kappa_{\mathrm{ind}}e^{-\bar s}\bigl(1+\frac12\mathrm{Var}_m(s)\bigr)
\)，
整理得到所示上界。
严格取等条件由 \(e^{-x}=1-x+\frac{x^2}{2}\) 仅在 \(x=0\) 处取等给出，从而等价于 \(x_j\equiv 0\)（按重数），即 \(\mathrm{Var}_m(s)=0\)。证毕。
\end{proof}

\begin{corollary}[可检测度函数：归一化峰值幅度是熵贡献的严格二次像]\label{cor:xi-detectability-amplitude-is-quadratic-image}
\leanverified{paper\_xi\_detectability\_amplitude\_is\_quadratic\_image}
对单原子剖面核
\(
H_{(\gamma,\delta,m)}(x)=m\,\frac{4\delta}{(x-\gamma)^2+(1+\delta)^2}
\)
定义归一化峰值幅度
\[
a:=\frac{H_{(\gamma,\delta,m)}(\gamma)}{m}=\frac{4\delta}{(1+\delta)^2}\in(0,1].
\]
令熵贡献
\(
p:=\frac{\delta}{1+\delta}\in(0,1)
\)，则成立严格代数同一
\begin{equation}\label{eq:xi-amplitude-as-quadratic-image-of-entropy}
\boxed{\;
a=4p(1-p)\; }.
\end{equation}
因此峰值幅度完全由熵贡献决定，并在 \(p=\tfrac12\)（等价 \(\delta=1\)）处达到最大值 \(a_{\max}=1\)。
反之，给定 \(a\in(0,1]\)，熵贡献仅有两支可能
\begin{equation}\label{eq:xi-entropy-from-amplitude-two-branches}
\boxed{\;
p=\frac12\bigl(1\pm\sqrt{1-a}\bigr)\; }.
\end{equation}
该二支分歧对应“仅凭幅度无法判定深度位于 \(\delta<1\) 还是 \(\delta>1\)”的固有二值不确定性；其消歧需借助相位/极点方向等解析信息（参见命题 \ref{prop:xi-lorentzian-rigidity-amplitude-halfwidth} 的分支讨论）。
\end{corollary}

\begin{proof}
由 \(p=\delta/(1+\delta)\) 得 \(1-p=1/(1+\delta)\)，故
\(
4p(1-p)=4\delta/(1+\delta)^2=a
\)，得到 \eqref{eq:xi-amplitude-as-quadratic-image-of-entropy}。
对 \eqref{eq:xi-entropy-from-amplitude-two-branches}，将 \(a=4p(1-p)\) 视为关于 \(p\) 的二次方程并解之即得；最大值与分支由抛物线开口向下与对称性给出。证毕。
\end{proof}

\begin{proposition}[粗粒化熵曲线的完全单调性：尺度轨道为 Stieltjes 指纹]\label{prop:xi-defect-entropy-scale-fingerprint-stieltjes}
引入尺度参数 \(t\ge 0\) 的 Poisson 粗粒化，把每个原子核的半宽由 \(1+\delta_j\) 改写为 \(t+1+\delta_j\)，并定义尺度熵
\begin{equation}\label{eq:xi-defect-entropy-scale-curve}
\boxed{\;
S_{\mathrm{def}}(t):=\sum_{j=1}^{\kappa}m_j\,\frac{\delta_j}{t+1+\delta_j}\;}
\end{equation}
（故 \(S_{\mathrm{def}}(0)=S_{\mathrm{def}}\) 且 \(S_{\mathrm{def}}(t)\downarrow 0\) 当 \(t\to\infty\)）。
则 \(t\mapsto S_{\mathrm{def}}(t)\) 为完全单调函数：对一切整数 \(n\ge 0\) 有
\begin{equation}\label{eq:xi-defect-entropy-scale-complete-monotone}
\boxed{\;
(-1)^n\frac{d^n}{dt^n}S_{\mathrm{def}}(t)
=n!\sum_{j=1}^{\kappa}m_j\,\frac{\delta_j}{(t+1+\delta_j)^{n+1}}
\ge 0\; }.
\end{equation}
因此 \(S_{\mathrm{def}}(t)\) 是一条可审计的 Stieltjes 型指纹曲线：在有限缺陷口径下，
其在复平面的有理延拓的极点集合与留数完全决定 \(\{(\delta_j,m_j)\}\)（忽略排列并合并同深度原子），
而高度参数 \(\{\gamma_j\}\) 在该标量尺度轨道中不出现。
\end{proposition}

\begin{proof}
对每一项 \(\frac{\delta_j}{t+1+\delta_j}\) 直接求导得到
\(
\frac{d^n}{dt^n}\frac{\delta_j}{t+1+\delta_j}
=(-1)^n n!\,\frac{\delta_j}{(t+1+\delta_j)^{n+1}}
\)，加权求和即得 \eqref{eq:xi-defect-entropy-scale-complete-monotone}，从而完全单调性成立。
\par
若进一步假设缺陷有限，则 \(S_{\mathrm{def}}(t)\) 为有理函数，
其极点位于 \(t=-1-\delta_j\)，留数为 \(m_j\delta_j\)。
因此由极点位置恢复 \(\{\delta_j\}\)，再由留数除以 \(\delta_j\) 恢复 \(\{m_j\}\)（合并同 \(\delta\) 情形下得到总重数）。
高度参数 \(\gamma_j\) 在 \eqref{eq:xi-defect-entropy-scale-curve} 中不出现，故不可能由该标量轨道辨识。证毕。
\end{proof}

\begin{proposition}[尺度熵流的 Riccati 夹逼：粗粒化导致的必然回收速率由 \texorpdfstring{$(\kappa_{\mathrm{ind}},\delta_{\min},\delta_{\max})$}{(kappa,delta_min,delta_max)} 封口]\label{prop:xi-defect-entropy-riccati-sandwich}
\leanverified{paper\_xi\_defect\_entropy\_riccati\_sandwich}
在命题 \ref{prop:xi-defect-entropy-scale-fingerprint-stieltjes} 的设定下，记
\[
\kappa_{\mathrm{ind}}:=\sum_{j=1}^{\kappa}m_j,
\qquad
\delta_{\min}:=\min_{1\le j\le\kappa}\delta_j,
\qquad
\delta_{\max}:=\max_{1\le j\le\kappa}\delta_j.
\]
则 \(S_{\mathrm{def}}(t)\) 可微且满足精确导数式
\[
S_{\mathrm{def}}'(t)=-\sum_{j=1}^{\kappa}m_j\,\frac{\delta_j}{(t+1+\delta_j)^2}.
\]
并且对一切 \(t\ge 0\) 成立双侧 Riccati 型夹逼
\begin{equation}\label{eq:xi-defect-entropy-riccati-sandwich}
\boxed{\;
-\frac{1}{\kappa_{\mathrm{ind}}\delta_{\min}}\,S_{\mathrm{def}}(t)^2
\ \le\
S_{\mathrm{def}}'(t)
\ \le\
-\frac{1}{\kappa_{\mathrm{ind}}\delta_{\max}}\,S_{\mathrm{def}}(t)^2.\;}
\end{equation}
等价地，
\begin{equation}\label{eq:xi-defect-entropy-inv-linear-bracket}
\boxed{\;
\frac{1}{S_{\mathrm{def}}(t)}-\frac{1}{S_{\mathrm{def}}(0)}
\ \in\
\Bigl[\frac{t}{\kappa_{\mathrm{ind}}\delta_{\max}},\ \frac{t}{\kappa_{\mathrm{ind}}\delta_{\min}}\Bigr].\;}
\end{equation}
因此尺度 \(t\) 的增长强制 \(S_{\mathrm{def}}(t)\) 以可证书的 \(t^{-1}\) 级衰减进入回收态；该衰减速率在制度层面仅依赖缺陷指数与深度区间，而与高度参数 \(\{\gamma_j\}\) 解耦。
\end{proposition}
\begin{proof}
令
\(
B_j(t):=m_j\frac{\delta_j}{(t+1+\delta_j)^2}>0
\)。
则
\[
-S_{\mathrm{def}}'(t)=\sum_{j=1}^{\kappa}B_j(t),
\qquad
S_{\mathrm{def}}(t)=\sum_{j=1}^{\kappa}(t+1+\delta_j)\,B_j(t),
\]
从而比值
\[
\frac{S_{\mathrm{def}}(t)}{-S_{\mathrm{def}}'(t)}
=\frac{\sum_j (t+1+\delta_j)B_j(t)}{\sum_j B_j(t)}
\]
是 \(\{t+1+\delta_j\}\) 的加权平均，故
\begin{equation}\label{eq:xi-defect-entropy-weighted-average-bracket}
t+1+\delta_{\min}
\ \le\
\frac{S_{\mathrm{def}}(t)}{-S_{\mathrm{def}}'(t)}
\ \le\
t+1+\delta_{\max}.
\end{equation}
另一方面，由单调性 \(\delta_{\min}\le \delta_j\le \delta_{\max}\) 与函数 \(\delta\mapsto \delta/(t+1+\delta)\) 递增，得
\[
\kappa_{\mathrm{ind}}\frac{\delta_{\min}}{t+1+\delta_{\min}}
\ \le\
S_{\mathrm{def}}(t)
\ \le\
\kappa_{\mathrm{ind}}\frac{\delta_{\max}}{t+1+\delta_{\max}}.
\]
将其分别与 \eqref{eq:xi-defect-entropy-weighted-average-bracket} 的左右端相乘得到
\[
\kappa_{\mathrm{ind}}\delta_{\min}
\ \le\
\frac{S_{\mathrm{def}}(t)^2}{-S_{\mathrm{def}}'(t)}
\ \le\
\kappa_{\mathrm{ind}}\delta_{\max},
\]
即
\(
\frac{1}{\kappa_{\mathrm{ind}}\delta_{\max}}
\le
(-S_{\mathrm{def}}'(t))/S_{\mathrm{def}}(t)^2
\le
\frac{1}{\kappa_{\mathrm{ind}}\delta_{\min}}
\)，等价于 \eqref{eq:xi-defect-entropy-riccati-sandwich}。
对恒等式
\(
\frac{d}{dt}\frac{1}{S_{\mathrm{def}}(t)}=-\frac{S_{\mathrm{def}}'(t)}{S_{\mathrm{def}}(t)^2}
\)
积分即得 \eqref{eq:xi-defect-entropy-inv-linear-bracket}。证毕。
\end{proof}

\endinput


\begin{theorem}[双 Herglotz 残数比定律：无需分离假设的 \texorpdfstring{$\delta$}{delta} 识别机制]\label{thm:xi-double-herglotz-residue-ratio}
在上述设定下，定义
\[
M_H(z):=\int_{\RR}\frac{H(x)}{x-z}\,\dd x,
\qquad
M_P(z):=-2\mathrm{i}\,\frac{d}{dz}\log S_{\mathrm{HP}}(z),
\qquad \Im z>0.
\]
则二者都可在 \(\CC\setminus\RR\) 上作有理延拓，并在下半平面极点
\[
z_j^-:=\gamma_j-\mathrm{i}(1+\delta_j)
\]
处具有同一极点集合；并且残数满足
\[
\operatorname*{Res}_{z=z_j^-} M_H(z)= -4\pi\,m_j\,\frac{\delta_j}{1+\delta_j},
\qquad
\operatorname*{Res}_{z=z_j^-} M_P(z)= 2\mathrm{i}\,m_j.
\]
因而成立严格残数比识别律
\[
\boxed{\;
\frac{\delta_j}{1+\delta_j}
\ =\
\frac{1}{2\pi}\left|\frac{\operatorname*{Res}_{z=z_j^-} M_H(z)}{\operatorname*{Res}_{z=z_j^-} M_P(z)}\right|\; }.
\]
该公式不依赖 \(\{\gamma_j\}\) 的最小分离：只要极点互异，则 \(\delta_j\) 在解析意义下可由残数直接识别；
分离仅影响数值稳定性（见推论 \ref{cor:xi-scan-collision-illposedness}）。
\end{theorem}

\begin{proof}
由 \eqref{eq:xi-comoving-scan-profile-H}，将单峰写成
\(
H_{(\gamma,\delta,m)}(x)=m\,\frac{4\delta}{(x-\gamma)^2+w^2}
\)
（\(w=1+\delta\)），直接用留数法计算其 Cauchy 变换得
\(
\int_{\RR}\frac{H_{(\gamma,\delta,m)}(x)}{x-z}\dd x
=-4\pi m\frac{\delta}{w}\cdot\frac{1}{z-(\gamma-\mathrm{i}w)}
\)，
故 \(M_H\) 的极点位于 \(z^-=\gamma-\mathrm{i}(1+\delta)\) 且残数如上。
\par
另一方面，
\(
\frac{d}{dz}\log S_{\mathrm{HP}}(z)
:=\sum_{j}m_j\left(\frac{1}{z-(\gamma_j+\mathrm{i}(1+\delta_j))}-\frac{1}{z-(\gamma_j-\mathrm{i}(1+\delta_j))}\right)
\)，
故在 \(z_j^-\) 处残数为 \(-m_j\)，从而 \(M_P\) 的残数为 \(2\mathrm{i}m_j\)。
代入即得残数比识别律。证毕。
\end{proof}

\begin{corollary}[单剖面全重建：共动密度 \(H\) 唯一决定点状缺陷多重集及其内函数编码]\label{cor:xi-comoving-profile-complete-reconstruction}
在有限点状缺陷口径下，给定共动扫描剖面 \(H\)（\eqref{eq:xi-comoving-scan-profile-H}），其 Cauchy 变换
\(
M_H(z)=\int_{\RR}\frac{H(x)}{x-z}\dd x
\)
在 \(\CC\setminus\RR\) 上作有理延拓，并在下半平面具有极点集合
\(
z_j^-=\gamma_j-\mathrm{i}(1+\delta_j)
\)
及残数闭式（定理 \ref{thm:xi-double-herglotz-residue-ratio}）：
\[
\operatorname*{Res}_{z=z_j^-} M_H(z)= -4\pi\,m_j\,\frac{\delta_j}{1+\delta_j}.
\]
因此 \(H\) 在解析意义下唯一确定多重集 \(\{(\gamma_j,\delta_j,m_j)\}_{j=1}^{\kappa}\)：
\(\gamma_j=\Re z_j^-\)、\(\delta_j=-\Im z_j^- -1\) 由极点位置给出，而
\begin{equation}\label{eq:xi-comoving-profile-recover-multiplicity-by-residue}
\boxed{\;
m_j
=
-\frac{1+\delta_j}{4\pi\,\delta_j}\operatorname*{Res}_{z=z_j^-} M_H(z)\ \in\ (0,\infty)\;}
\end{equation}
由残数与 \(\delta_j\) 唯一恢复。
从而上半平面内函数编码
\[
S_{\mathrm{HP}}(z)=\prod_{j=1}^{\kappa}\left(\frac{z-\gamma_j-\mathrm{i}(1+\delta_j)}{z-\gamma_j+\mathrm{i}(1+\delta_j)}\right)^{m_j}
\]
亦由 \(H\) 唯一确定（以归一化 \(S_{\mathrm{HP}}(\infty)=1\) 固定单位模常数）。
在 Cayley 逆像下，盘内点
\(
a_j=\mathcal{C}(\gamma_j+\mathrm{i}\delta_j)\in\mathbb{D}
\)
由 \((\gamma_j,\delta_j)\) 唯一确定；故 \(H\) 是“点状缺陷”意义下的完整边界不变量。
\end{corollary}

\begin{corollary}[边界重建的整数化判别：由残数整除性锁定重数]\label{cor:xi-comoving-profile-integrality-test}
在推论 \ref{cor:xi-comoving-profile-complete-reconstruction} 的设定下，若仅给定 \(H\) 并要求其来自某个\emph{整数重数}点状缺陷机制（即 \(m_j\in\NN\)），
则必要且充分条件为：\(M_H\) 的极点集合满足 \(\delta_j=-\Im z_j^- -1>0\)，并且对每个极点 \(z_j^-\) 都有
\[
\boxed{\;
-\frac{1+\delta_j}{4\pi\,\delta_j}\operatorname*{Res}_{z=z_j^-} M_H(z)\ \in\ \NN.\;}
\]
该条件为“边界数据是否可由合法点状缺陷产生”的纯解析可判定准则：极点位置给出连续深度，而残数的整除性给出离散重数。
\end{corollary}

\begin{proof}
由 \eqref{eq:xi-comoving-profile-recover-multiplicity-by-residue}，\(m_j\) 的闭式恢复是必要的；其取值于 \(\NN\) 亦显然充分。证毕。
\end{proof}

\begin{proposition}[无穷远剖面矩展开：扫描尾部系数给出轴外参数的可审计矩（50）]\label{prop:xi-scan-far-field-moment-expansion}
当 \(|x|\to\infty\) 时，共动扫描剖面 \(H(x)\) 具有渐近展开
\[
H(x)\sim \sum_{n\ge 0}\frac{M_n}{x^{n+2}},
\]
其中系数 \(M_n\) 为缺陷参数的显式代数矩。
前三项为
\[
\boxed{\;
M_0=4\sum_{j=1}^{\kappa}m_j\delta_j,\qquad
M_1=8\sum_{j=1}^{\kappa}m_j\delta_j\,\gamma_j,\qquad
M_2=4\sum_{j=1}^{\kappa}m_j\delta_j\bigl(3\gamma_j^2-(1+\delta_j)^2\bigr)\;}
\]
并且一般地 \(M_n\) 可写为 \(\sum_j m_j\delta_j\) 与 \(\sum_j m_j\delta_j\,\gamma_j^k(1+\delta_j)^\ell\)（\(k+\ell\le n\)）的线性组合。
因此若能在远场区间上以足够精度估计 \(M_0,M_1,\dots,M_{N}\)，即可得到一列可审计的“矩链证书”，把二维离线参数转写为标量代数不变量。
\end{proposition}

\begin{proof}
对单峰作 \(|x|\to\infty\) 展开：
\[
\frac{4\delta}{(x-\gamma)^2+(1+\delta)^2}
:=\frac{4\delta}{x^2}\cdot\frac{1}{1-\frac{2\gamma}{x}+\frac{\gamma^2+(1+\delta)^2}{x^2}},
\]
再对括号内作幂级数展开并逐项求和即可得到 \(x^{-(n+2)}\) 系数为关于 \(\gamma\) 与 \(1+\delta\) 的多项式。
汇总各峰并按重数加权即得结论；\(M_0,M_1,M_2\) 由前三项直接读出。证毕。
\end{proof}

\begin{proposition}[远场矩—极点连锁：尾系数满足最小阶 \(2\kappa\) 线性递推]\label{prop:xi-far-field-moment-recursion-order-2kappa}
\leanverified{paper\_xi\_far\_field\_moment\_recursion\_order\_2kappa}
在有限点状缺陷口径下，\(H\) 为实有理函数且可写为 \(H(x)=P(x)/Q(x)\)，其中
\[
Q(x)=\prod_{j=1}^{\kappa}\bigl((x-\gamma_j)^2+(1+\delta_j)^2\bigr),
\qquad \deg Q=2\kappa,\ \deg P\le 2\kappa-2.
\]
令远场展开 \(H(x)\sim\sum_{n\ge 0}\frac{M_n}{x^{n+2}}\)（命题 \ref{prop:xi-scan-far-field-moment-expansion}），并定义反转多项式
\begin{equation}\label{eq:xi-Qsharp-from-Q}
\boxed{\;
Q^\sharp(y):=y^{2\kappa}Q(1/y)=\sum_{r=0}^{2\kappa}q_r\,y^{r}\qquad(q_0=1).\;}
\end{equation}
则尾系数满足严格递推
\begin{equation}\label{eq:xi-far-field-moment-recursion}
\boxed{\;
\sum_{r=0}^{2\kappa} q_r\,M_{n+r}=0\qquad(\forall n\ge 0).\;}
\end{equation}
并且当 \(P,Q\) 既约（等价于 \(H\) 的极点无消去）时，\eqref{eq:xi-far-field-moment-recursion} 的最小递推阶数恰为 \(2\kappa\)；
因此 \(\kappa\) 等于尾矩序列的最小递推阶数的一半，而 \(Q\)（从而全部 \((\gamma_j,\delta_j)\)）可由有限段 \(\{M_n\}\) 通过 Pad\'e/消元（或等价的线性代数）完全恢复。
\end{proposition}

\begin{proof}
令 \(y=1/x\)，并设
\(
F(y):=y^{-2}H(1/y)=\sum_{n\ge 0}M_ny^n
\)。
由 \(H=P/Q\) 得
\(
y^{2\kappa-2}P(1/y)=Q^\sharp(y)\,F(y)
\)，其中左端为次数不超过 \(2\kappa-2\) 的多项式。
故 \(Q^\sharp(y)F(y)\) 的幂级数在 \(y^{2\kappa-1}\) 及以上系数全为零，即得 \eqref{eq:xi-far-field-moment-recursion}。
若 \(H\) 的最小分母次数为 \(2\kappa\)，则 \(Q^\sharp\) 亦为最小；从而递推阶数不可再降。证毕。
\end{proof}

\begin{proposition}[Poisson 粗粒化的无毛发极限：大尺度 \texorpdfstring{$t$}{t} 下仅剩有限矩不变量]\label{prop:xi-poisson-coarsegraining-nohair}
定义缺陷测度在高度方向的低阶矩
\[
M_0^{\mathrm{def}}:=\sum_{j=1}^{\kappa}m_j\delta_j,\qquad
M_1^{\mathrm{def}}:=\sum_{j=1}^{\kappa}m_j\delta_j\,\gamma_j,\qquad
\bar\gamma:=\frac{M_1^{\mathrm{def}}}{M_0^{\mathrm{def}}}\quad(M_0^{\mathrm{def}}>0).
\]
令 Poisson 半群粗粒化剖面
\[
G_t(x)
:=\frac{1}{\pi}\sum_{j=1}^{\kappa}m_j\delta_j\,\frac{t}{(x-\gamma_j)^2+t^2},
\qquad t>0.
\]
则当 \(t\to\infty\) 时成立一致渐近式
\[
\boxed{\;
G_t(x)
=\frac{M_0^{\mathrm{def}}}{\pi}\,\frac{t}{(x-\bar\gamma)^2+t^2}
\ +\ O(t^{-2}),\;}
\]
其中余项 \(O(t^{-2})\) 在 \(x\in\RR\) 上一致，并由二阶中心矩
\(\sum_j m_j\delta_j(\gamma_j-\bar\gamma)^2\)
控制。
因此在大尺度粗粒化下，全部点状缺陷在制度可见域中不可避免地退化为单个“有效 Lorentzian”，其可区分参数仅为总质量 \(M_0^{\mathrm{def}}\) 与重心 \(\bar\gamma\)；更高阶几何信息被严格压制到 \(t^{-2}\) 级别。
\end{proposition}

\begin{proof}
记 \(w_j:=m_j\delta_j\) 并令 \(f_t(x;\gamma):=\frac{t}{(x-\gamma)^2+t^2}\)。
对 \(\gamma\) 在 \(\bar\gamma\) 处作 Taylor 展开：
\[
f_t(x;\gamma_j)
:=f_t(x;\bar\gamma)
+f_t'(x;\bar\gamma)(\gamma_j-\bar\gamma)
+\frac12 f_t''(x;\xi_j)(\gamma_j-\bar\gamma)^2
\]
其中 \(\xi_j\) 介于 \(\gamma_j\) 与 \(\bar\gamma\) 之间。
加权求和并用 \(\sum_j w_j(\gamma_j-\bar\gamma)=0\) 消去一阶项，得
\[
\sum_j w_j f_t(x;\gamma_j)
:=M_0^{\mathrm{def}} f_t(x;\bar\gamma)
+\frac12\sum_j w_j f_t''(x;\xi_j)(\gamma_j-\bar\gamma)^2.
\]
而对任意 \(t>0\) 有一致界 \(\sup_{x,\gamma}|f_t''(x;\gamma)|\lesssim t^{-3}\)，故余项一致为 \(O(t^{-3})\)，从而亦为 \(O(t^{-2})\)；
常数由二阶中心矩控制。再乘以 \(1/\pi\) 即得结论。证毕。
\end{proof}

\leanverified{paper\_xi\_toeplitz\_vs\_scan\_linear\_lowerbound}
\begin{proposition}[两类证书的统一线性下界：Toeplitz 阶数与扫描覆盖数]\label{prop:xi-toeplitz-vs-scan-linear-lowerbound}
固定偏移下界 \(\delta\ge\delta_0>0\) 与高度窗 \(|\gamma|\le T\)。
若要求一种制度内证书对任意可能的缺陷位置给出一致可检测性，则：
\begin{enumerate}
  \item \textbf{（扫描覆盖的线性门槛）}由于每一峰的半宽为 \(1+\delta_j\ge 1+\delta_0\)，任一有限扫描集合 \(\{x_0^{(k)}\}_{k=1}^{K}\subset[-T,T]\) 若要保证对任意 \(\gamma\in[-T,T]\) 都存在某个 \(k\) 使 \(|\gamma-x_0^{(k)}|\le 1+\delta_0\)，则必有
  \[
  \boxed{\;K\ \gtrsim\ \frac{T}{1+\delta_0}\;}
  \]
  （覆盖数至少线性增长）。
  \item \textbf{（Toeplitz 阶数的线性必要性）}在固定图表（不外置共动地址）的 Toeplitz--PSD 口径下，统一排除高度窗内潜在离切片极点的矩阶满足更强的必要下界（推论 \ref{cor:xi-toeplitz-height-phi-calibration}）；特别地 \(N=\Omega(T)\) 为必要条件。
\end{enumerate}
因此两类证书在资源下界意义下至少同为线性增长；其差异不在于是否存在下界，而在于是否允许以“共动地址”把端点红移压缩外置为覆盖资源。
\end{proposition}

\begin{proof}
第 (1) 条为区间覆盖的直接计数：长度 \(2T\) 的区间用半径 \((1+\delta_0)\) 的球覆盖所需个数为 \(\Omega(T/(1+\delta_0))\)。
第 (2) 条由推论 \ref{cor:xi-toeplitz-height-phi-calibration} 立即推出。证毕。
\end{proof}
